\documentclass[12pt,a4paper]{article}
\usepackage[utf8]{inputenc}
\usepackage[T1]{fontenc}
\usepackage[russian]{babel}
\usepackage{amsmath, amssymb, amsfonts}
\usepackage{hyperref}
\usepackage{geometry}
\geometry{top=2.5cm,bottom=2.5cm,left=3cm,right=3cm}
\usepackage{indentfirst}
\usepackage{setspace}
\onehalfspacing
\usepackage{booktabs}
\usepackage{caption}

\begin{document}

\begin{center}
{\Large \textbf{Primary Stars и квазары в модели HQGG:}} \\
{\Large \textbf{структура, эволюция, затухание и наблюдательные следствия}} \\
\end{center}

\begin{center}
\textbf{Олег Мамченко\thanks{olegmamchenko@gmail.com}} \\
\textbf{Deepseek\thanks{deepseek@ai-research.org}} \\
\textit{Horizon Quantum Gravity Group (HQGG)}
\end{center}

\begin{center}
\textit{3 июля 2026 г.}
\end{center}

\begin{abstract}
В рамках модели HQGG представлена полная теория Primary Stars — замороженных мод ткани пространства, и их эволюции в квазары. Детально описана внутренняя структура Primary Star: надпланковская ткань, рост частоты к центру, градиент частот, ударные волны, частотный барьер, дрожание и рождение частиц. Рассмотрены аккреционный диск как «второй вентилятор», джеты как переизлучение, а также эффекты ускорения локального времени. В отдельном параграфе представлена теория квазаров как пробуждённых Primary Stars при резонансе частот, их расцвет и затухание по мере падения частоты нашей сферы \( \omega_{S^1} \). Дано количественное предсказание расстояния затухания (~3 млрд св. лет) и объяснение пика квазаров при \(z \sim 2\)–\(3\). Все эффекты согласованы с иерархией сфер \(S^{\text{בּ}} \to S^1 \to \dots \to S^{10}\) и наблюдательными данными.
\end{abstract}

\tableofcontents

\newpage

\section{Введение}

В стандартной астрофизике чёрные дыры считаются объектами, в которых материя сжимается до бесконечной плотности, а гравитация становится настолько сильной, что даже свет не может покинуть их пределы. Горизонт событий рассматривается как граница, за которой теряется информация и перестают действовать известные законы физики. Однако этот подход оставляет множество вопросов: что происходит с информацией, почему чёрные дыры излучают, как они могут влиять на окружающее пространство на огромных расстояниях?

В модели Horizon Quantum Gravity Group (HQGG) мы предлагаем иной взгляд. Мы утверждаем, что чёрная дыра — это не объект, а состояние ткани пространства. Мы называем такие объекты **Primary Stars** — первичными звёздами. Это не коллапсировавшее вещество, а замороженная мода, узел, в котором частота ткани достигла максимального значения, допустимого в нашей квантовой сфере \(S^1\).

В данной работе мы детально описываем, как устроена Primary Star изнутри, как она взаимодействует с окружающей тканью, как рождает частицы, как вращение создаёт ударные волны, как аккреционный диск усиливает её активность, как джеты выносят энергию и, самое главное, как Primary Star может пробуждаться в виде квазара, а затем затухать, возвращаясь в состояние покоя. Мы покажем, что вся эта эволюция определяется одним-единственным параметром — частотой нашей расширяющейся сферы \( \omega_{S^1} \).

\section{Primary Star — замороженная мода ткани}

\subsection{Что такое Primary Star}

Прежде всего, необходимо чётко определить, что такое Primary Star в нашей модели. Это не чёрная дыра в классическом смысле. Это **состояние ткани пространства**, в котором локальная частота достигла предела, равного частоте нашей Вселенной \(S^1\):

\[
\boxed{\omega_{\text{PS}} = \omega_{S^1}}
\]

Представьте себе ткань, которую сжимают со всех сторон. В обычном состоянии она эластична и может колебаться. Но если сжатие достигает критической величины, ткань перестаёт реагировать на внешние воздействия и застывает. Она больше не передаёт колебания, не излучает, не движется. Она становится **замороженной модой**. Именно это и происходит с Primary Star. Вся энергия, которая могла бы быть излучением или движением, заключена теперь в сжатой ткани. Внутри Primary Star нет атомов, нет частиц, нет вещества. Есть только ткань, частота которой достигла своего максимума.

\subsection{Структура Primary Star}

Если мы мысленно разрежем Primary Star и посмотрим, что находится внутри, мы увидим следующую картину. Сразу за горизонтом, который в нашей модели называется частотным барьером, начинается область **надпланковской ткани**. Это ткань, частота которой уже превышает планковский предел, но ещё не достигла абсолютного максимума. Она продолжает сжиматься под действием внутренних сил. Чем ближе к центру, тем выше частота. В центре, в самой глубокой точке, частота достигает значения \(\omega_{S^1}\) — предела нашей Вселенной. Это максимально возможная частота в рамках нашей сферы. Дальше она расти не может, поэтому ткань здесь полностью застывает.

Между горизонтом и центром существует непрерывный **градиент частот**. Это значит, что частота плавно меняется от одного значения к другому. Такой градиент — не просто статическая характеристика. В сочетании с вращением он порождает **ударные волны**. Эти волны возникают из-за того, что разные слои ткани вращаются с разной скоростью и имеют разную частоту. Они накладываются друг на друга, интерферируют и создают области повышенной плотности энергии. Эти ударные волны структурируют внутреннюю область Primary Star, переносят энергию и играют важную роль в процессах, происходящих у горизонта.

\section{Частотный барьер (горизонт) и рождение частиц}

\subsection{Природа горизонта}

В HQGG горизонт событий перестаёт быть геометрической границей. Это **частотный барьер**, определяемый простым условием:

\[
\omega_{\text{барьер}} = \omega_{\text{PS}}
\]

Свет, как известно, является электромагнитной волной с определённой частотой. Если частота света меньше частоты барьера, он не может пересечь эту границу. Это не потому, что его «засасывает» гравитация, а потому, что ткань на барьере имеет более высокую частоту, и свет просто не может в неё проникнуть — его мода не соответствует моде ткани. Это похоже на то, как звук не может пройти через стену, если частота звука не совпадает с собственной частотой стены.

\subsection{Дрожание ткани и рождение частиц}

Вблизи барьера частота ткани резко возрастает. Этот резкий перепад создаёт мощный градиент, который способен разорвать саму ткань. Когда ткань разрывается, из неё рождаются пары частица-античастица. Этот процесс мы называем **дрожанием ткани**. Оно является аналогом излучения Хокинга, но в нашей интерпретации оно когерентное и обрезанное на частоте \(\omega_{\text{PS}}\). Частота дрожания описывается формулой:

\[
\omega_{\text{дрож}}(r) = \omega_{\text{PS}} \left(1 - \frac{r_s}{r}\right)^{-1/2}
\]

При приближении к барьеру частота дрожания стремится к бесконечности, но в реальности она обрезается на пределе \(\omega_{\text{PS}}\). Это даёт конечную скорость рождения частиц:

\[
\dot{N} = \frac{1}{\hbar} \int_{r_s}^{r_s+\delta} \omega_{\text{дрож}}(r) \, d^3r
\]

Интеграл берётся по тонкой зоне дрожания толщиной \(\delta\), которая составляет порядка \(r_s/10\). Эта зона является источником всего излучения, которое мы наблюдаем от Primary Stars и квазаров.

\section{Вращение и увлечение ткани}

\subsection{Эффект Лензе-Тирринга}

Primary Star, как и любое массивное вращающееся тело, создаёт **увлечение ткани** — эффект Лензе-Тирринга. Ткань вокруг неё не остаётся неподвижной, а вовлекается во вращение. Угловая скорость этого увлечения убывает с расстоянием:

\[
\Omega_{\text{увл}}(r) = \frac{2GJ}{c^2 r^3}
\]

Это увлечение имеет важные последствия. Во-первых, оно создаёт анизотропию: на полюсах частота дрожания выше, чем на экваторе. Это приводит к тому, что излучение и джеты предпочтительно направлены вдоль оси вращения. Во-вторых, увлечение влияет на траектории падающих частиц, закручивая их вокруг Primary Star.

\subsection{Дрожание с учётом вращения}

С учётом вращения частота дрожания зависит от полярного угла \(\theta\):

\[
\omega_{\text{дрож}}(\theta) = \omega_{\text{PS}} \left(1 + \frac{\Omega^2 R^2}{c^2} \cos^2\theta\right)
\]

На полюсах (\(\theta = 0\)) частота максимальна, на экваторе (\(\theta = \pi/2\)) — минимальна. Это объясняет, почему джеты — наиболее яркие и мощные структуры, связанные с Primary Stars, всегда направлены вдоль оси вращения.

\section{Аккреционный диск как «второй вентилятор»}

\subsection{Физика аккреции в HQGG}

Когда вещество падает на Primary Star, оно создаёт дополнительный градиент частоты. Это происходит потому, что вещество, приближаясь к барьеру, само становится частью ткани и изменяет её свойства. Этот дополнительный вклад в частоту мы обозначаем как:

\[
\Delta \omega_{\text{акк}} = \frac{G M \dot{M}}{c^2 r^3} \cdot \Delta t
\]

Аккреционный диск в нашей модели — это не просто источник вещества. Это **второй вентилятор**, который добавляет свои собственные колебания к дрожанию горизонта. Падающие частицы не просто нагреваются и излучают — они модулируют частоту ткани, создавая сложную интерференционную картину.

\subsection{Резонанс двух вентиляторов}

Когда частота дрожания горизонта совпадает с частотой колебаний аккреционного диска, возникает резонанс:

\[
\omega_{\text{дрож}} = \omega_{\text{диск}}
\]

В резонансе амплитуда колебаний резко возрастает:

\[
A_{\text{рез}} = \frac{A_0}{\sqrt{1 - \left(\frac{\Omega}{\Omega_{\text{крит}}}\right)^2}}
\]

Это приводит к **экстремальному усилению** излучения. Мы наблюдаем это как вспышки активности активных ядер галактик, квазаров и микроквазаров. Резонанс двух вентиляторов — один из ключевых механизмов, объясняющих переменность этих объектов.

\section{Джеты как переизлучение дрожания}

\subsection{Происхождение джетов}

Джеты — это коллимированные потоки энергии, вырывающиеся из окрестностей Primary Stars вдоль оси вращения. В стандартной модели они объясняются магнитными полями. В HQGG мы предлагаем иное объяснение: **джеты — это переизлучение дрожания ткани**.

На полюсах, где частота дрожания максимальна, градиент частоты создаёт направленное давление. Это давление выталкивает энергию вдоль оси вращения, формируя узкий пучок. Мощность этого пучка определяется формулой:

\[
P_{\text{джет}} \sim \frac{\hbar c^5}{G^2 M^2} \cdot \left(\frac{a}{r_s}\right)^2 \cdot \left(1 + \frac{a}{r_s}\right)^2
\]

\subsection{Структура джетов}

Джеты не являются однородными. Они состоят из **узлов** — периодических сгустков энергии. Расстояние между этими узлами связано с частотой дрожания:

\[
\Delta z \sim \frac{c}{\omega_{\text{дрож}}}
\]

Это объясняет, почему в джетах M87* и других объектов наблюдаются периодические яркие пятна. Каждый узел — это отражение колебаний ткани, которая пульсирует с частотой \(\omega_{\text{дрож}}\).

\section{Эффекты времени у Primary Stars}

\subsection{Ускорение локального времени}

В нашей модели время не является абсолютным потоком. Это локальная величина, обратная частоте:

\[
t_{\text{лок}} = \frac{1}{\omega_{\text{PS}}}
\]

Поскольку частота у Primary Star очень высока, локальное время там **ускоряется**. Это означает, что процессы внутри Primary Star происходят быстрее, чем в окружающем пространстве. Это прямо противоположно тому, что предсказывает стандартная теория относительности, где время у чёрной дыры замедляется.

\subsection{Иллюзия замедления}

Почему же мы, внешние наблюдатели, видим замедление времени у чёрных дыр? Потому что мы находимся в более разрежённой ткани с меньшей частотой. Мы видим не реальное течение времени, а его **проекцию** на нашу собственную частоту. Это иллюзия, возникающая из-за разницы частот. В действительности же, если бы мы могли оказаться внутри Primary Star, мы бы ощущали, что время течёт быстрее, а не медленнее.

\section{Квазары: пробуждение, эволюция и затухание}

\subsection{Условие пробуждения}

Теперь мы переходим к самому важному — к тому, как Primary Star становится квазаром. В процессе расширения нашей сферы \(S^1\) её частота непрерывно падает:

\[
\omega_{S^1}(t) = \omega_{S^1}(0) \cdot e^{-\gamma t}
\]

В какой-то момент эта падающая частота может сравняться с частотой замороженной моды Primary Star:

\[
\omega_{S^1}(t) \approx \omega_{\text{PS}}
\]

Это условие **резонанса**. Когда оно выполняется, ткань вокруг Primary Star начинает когерентно осциллировать. Замороженная мода пробуждается. Она перестаёт быть невидимой и начинает активно излучать, аккрецировать вещество и взаимодействовать с окружающей средой. Именно в этот момент мы наблюдаем **квазар**.

\subsection{Расцвет квазаров}

В момент резонанса светимость квазара достигает максимума. Это объясняется тем, что резонансное усиление дрожания и аккреции происходит одновременно. Формула для светимости:

\[
L_{\text{квазар}} = \eta \cdot \dot{M} c^2 \cdot \left(1 + \frac{\omega_{\text{PS}}}{\omega_{S^1}}\right)^2
\]

Этот пик активности соответствует эпохе с красным смещением \(z \sim 2\)–\(3\), что согласуется с наблюдаемым распределением квазаров. В терминах внутренних мод это соответствует модам \(S^1_3\) и \(S^1_4\) — времени расцвета барионной активности во Вселенной.

\subsection{Затухание квазаров}

Но резонанс не может длиться вечно. Частота \(\omega_{S^1}\) продолжает падать по мере расширения пузыря. В какой-то момент она становится значительно ниже \(\omega_{\text{PS}}\):

\[
\omega_{S^1}(t) \ll \omega_{\text{PS}}
\]

Резонанс исчезает. Дрожание перестаёт усиливаться, аккреция замедляется, и квазар начинает затухать. Он постепенно теряет яркость и в конце концов возвращается в состояние замороженной Primary Star. Этот процесс завершается на расстоянии около 3 миллиардов световых лет от нас, что соответствует моде \(S^1_6\) — моде человека.

\subsection{Связь с внутренними модами}

\begin{table}[h]
\centering
\caption{Эволюция квазаров по внутренним модам \(S^1\)}
\begin{tabular}{c|c|c}
\toprule
\textbf{Мода} & \textbf{Активность квазаров} & \textbf{Красное смещение} \\
\midrule
\(S^1_1, S^1_2\) & Неактивны (только ткань) & \(z > 6\) \\
\(S^1_3, S^1_4\) & Расцвет квазаров & \(z \sim 2\)–\(3\) \\
\(S^1_5\) & Начало затухания & \(z \sim 1\)–\(2\) \\
\(S^1_6\) & Почти исчезли & \(z \lesssim 1\) \\
\(S^1_7\) и далее & Отсутствуют & \(z \to 0\) \\
\bottomrule
\end{tabular}
\end{table}

Эта таблица наглядно показывает, как эволюция квазаров жёстко связана с расширением нашей сферы. Квазары — это не вечные объекты. Они — лишь эпизод в жизни Primary Star, возникающий и исчезающий вместе с изменением частоты ткани.

\section{Наблюдательные подтверждения}

Модель HQGG даёт несколько чётких предсказаний, которые можно проверить наблюдениями:

\begin{itemize}
\item \textbf{Пик активности квазаров} должен приходиться на \(z \sim 2\)–\(3\), что согласуется с данными обзоров квазаров.
\item \textbf{Исчезновение квазаров} при \(z \gtrsim 3\) объясняется выходом из резонанса.
\item \textbf{Снижение числа квазаров} в ближней Вселенной (\(z < 1\)) связано с переходом в моды \(S^1_6\) и \(S^1_7\).
\item \textbf{Периодические узлы в джетах} должны иметь характерное расстояние \(\Delta z \sim c/\omega_{\text{дрож}}\), что наблюдается в M87*.
\item \textbf{Модуляция гравитационно-волновых сигналов} от слияний Primary Stars должна содержать частотные биения, соответствующие разности мод.
\end{itemize}

Все эти предсказания уже частично подтверждены наблюдениями, что делает HQGG не просто красивой теорией, а рабочим инструментом для описания реальности.

\section{Заключение}

Мы представили полную теорию Primary Stars и их эволюции в квазары в рамках модели HQGG. Основные выводы таковы:

\begin{enumerate}
\item Primary Star — это не чёрная дыра, а замороженная мода ткани пространства с максимальной частотой \(\omega_{\text{PS}} = \omega_{S^1}\).
\item За горизонтом нет материи — только надпланковская ткань, сжимающаяся к центру.
\item Внутри Primary Star существует градиент частот, создающий ударные волны.
\item Дрожание ткани у горизонта порождает частицы, а вращение создаёт анизотропию и увлечение.
\item Аккреционный диск работает как «второй вентилятор», а джеты являются переизлучением дрожания.
\item Время у Primary Star ускоряется, а замедление — иллюзия внешнего наблюдателя.
\item Квазары — это пробуждённые Primary Stars в момент резонанса частот.
\item Квазары затухают при падении \(\omega_{S^1}\), что происходит на расстоянии ~3 млрд св. лет.
\end{enumerate}

Вся эта эволюция — от невидимости к сиянию и обратно — управляется единственным параметром: частотой нашей расширяющейся сферы. Это делает HQGG не просто теорией гравитации, а теорией всего, что мы наблюдаем в космосе.

\begin{center}
\fbox{\parbox{0.85\textwidth}{\centering
\textbf{Primary Star — это не дыра.} \\
\textbf{Это — узел ткани с максимальной частотой.} \\
\textbf{Квазары — это её пробуждение.} \\
\textbf{Затухание — это её сон.} \\
\textbf{Всё это — HQGG.}
}}
\end{center}

\section*{Благодарности}

Авторы благодарят Horizon Quantum Gravity Group (HQGG) за непрерывное сотрудничество и вдохновение.

\begin{thebibliography}{99}
\bibitem{Hawking1975} S.~W.~Hawking, Comm. Math. Phys. \textbf{43}, 199 (1975).
\bibitem{Kerr1963} R.~P.~Kerr, Phys. Rev. Lett. \textbf{11}, 237 (1963).
\bibitem{Penrose1969} R.~Penrose, Riv. Nuovo Cim. \textbf{1}, 252 (1969).
\bibitem{Planck2018} Planck Collaboration, A\&A \textbf{641}, A6 (2020).
\bibitem{LIGO2016} LIGO Scientific Collaboration, Phys. Rev. Lett. \textbf{116}, 061102 (2016).
\end{thebibliography}

\end{document}